\chapter{Instrukcja uruchomienia projektu}
\label{ch:instrukcja}

Niniejszy rozdział zawiera kompletny przewodnik pozwalający na pobranie, skonfigurowanie i poprawne uruchomienie aplikacji Rock Gym w środowisku deweloperskim lub testowym.

\section{Wymagania systemowe i środowisko programistyczne}
Do bezproblemowego uruchomienia i skompilowania projektu wymagane jest spełnienie następujących warunków:
\begin{itemize}
    \item \textbf{System operacyjny:} Windows 11 64-bit (lub nowszy zgodny z platformą WPF).
    \item \textbf{Procesor:} klasy AMD Ryzen 5 5500U (lub odpowiednik o podobnej bądź wyższej wydajności).
    \item \textbf{Pamięć RAM:} minimum 8 GB.
    \item \textbf{Karta graficzna:} zintegrowana (brak szczególnych wymagań co do akceleratora).
    \item \textbf{Środowisko bazy danych:} silnik MySQL w wersji 5.2.1 lub wyższej.
    \item \textbf{Środowisko uruchomieniowe i programistyczne:} Microsoft Visual Studio (zalecana wersja 2026), zainstalowana platforma .NET w wersji 10.0 oraz biblioteka EntityFrameworkCore 9.0.0.
\end{itemize}

\section{Pobranie projektu i konfiguracja bazy danych}
Całość kodu źródłowego, wraz z kopią zapasową bazy danych niezbędną do odtworzenia środowiska, znajduje się w udostępnionym, publicznym repozytorium w serwisie GitHub:
\begin{center}
    \url{https://github.com/BartlomiejKufel/rock-gym-personnel}
\end{center}

\subsection{Odtworzenie bazy danych z backupu}
Wraz z repozytorium udostępniono plik powiązany ze schematem bazy danych gotowy do zaimportowania. Aby oprogramowanie działało poprawnie, należy zrealizować następujące kroki:
\begin{enumerate}
    \item Uruchomić lokalny serwer MySQL.
    \item Za pomocą klienta (np. phpMyAdmin) utworzyć nową pustą bazę danych i zaimportować do niej dołączony do projektu plik \texttt{rock\_gym.sql}. Proces ten zbuduje od podstaw strukturę wszystkich wymaganych tabel oraz zasili je wstępnymi danymi testowymi.
    \item Otworzyć rozwiązanie \textit{RockGym} w programie Visual Studio.
    \item Wyszukać lokalizację definicji połączenia (\textit{ConnectionString}) która znajduje się w module \texttt{RockGymContext.cs} i odpowiednio podmienić parametry dostępowe (adres serwera localhost, port, nazwa utworzonej bazy, domyślny użytkownik i hasło serwera).
\end{enumerate}

\section{Uruchomienie aplikacji i skanera biometrycznego}
Po poprawnym skonfigurowaniu \textit{ConnectionString}, projekt jest gotowy do uruchomienia.
\begin{enumerate}
    \item W środowisku Visual Studio należy przebudować projekt i rozpocząć debugowanie (np. korzystając ze skrótu \texttt{F5}).
    \item Po uruchomieniu powinien ukazać się standardowy ekran logowania dla personelu.
    \item W pobranym archiwum projektu umieszczono dodatkowy folder \texttt{fingerprints} zawierający graficzne próbki odcisków palców. Symulacja weryfikacji i autoryzacji odbywa się poprzez wczytanie jednego z dostępnych tam obrazów prosto do dedykowanego okna w aplikacji Rock Gym.
\end{enumerate}

\section{Dane logowania do kont testowych}
W zaimplementowanej testowej bazie danych przygotowano kilka kont dla pracowników, pozwalające na bezproblemowe przetestowanie funkcjonalności systemu bez tracenia czasu na ręczną dodawanie użytkowników. Poniżej wypisano dostępne domyślne dane autoryzacyjne:
\begin{itemize}
    \item \textbf{Konto administratora / menedżera (pełen dostęp):}
          \begin{itemize}
              \item Login: \texttt{admin}
              \item Hasło: \texttt{admin}
          \end{itemize}
          \vspace{0.3cm}
    \item \textbf{Konto pracownicze (dostęp ograniczony do operacji podstawowych):}
          \begin{itemize}
              \item Login: \texttt{pracownik}
              \item Hasło: \texttt{pracownik}
          \end{itemize}
\end{itemize}
